\problem{Магниттік тізбектер мен трансформатор}{10.0}

Ом заңы мектеп физикасынан бәріне мәлім. Бұл есепте сіз магнетизм тұрғысынан электр тізбектеріне ұқсас заңдар қолданылатын \textit{магниттік тізбектер} ұғымымен танысасыз және олардың көмегімен электротехникада қолданылатын кейбір құрылғыларды талдайсыз.

\pic5r{0pt}{0.2}{figures/3.1.png}

Сізге мәлім болған магниттік индукция $\vec B$ векторымен қатар, \textit{магнит өрісінің кернеулігі} $\vec H$ векторын енгізу пайдалы. Магнит өрісінің кернеулігі үшін Ампер теоремасы бойынша, $C$ контуры бойымен $\vec H$ векторының циркуляциясы тұйық контур арқылы өтетін \textit{қорытқы} өткізгіштік ток күшіне тең:
$$
\oint_C\vec H\cdot d\vec l = I.
$$
Бұл есепте барлық магнетиктер $\vec B = \mu\mu_0\vec H$ сызықтық заңына бағынады деп есептеңіз, мұндағы $\mu$ --- заттың магниттік өтімділігі, ал $\mu_0=4\pi\cdot10{-7}~\unit{\H\per\m}$ --- магниттік тұрақты. Есепте магнетик ретінде темір алынады, оның өтімділігі $\mu=2000$ тұрақты, ал ферромагниттік қанығу және гистерезис құбылыстарын ескермеуге болады.

\part{Жаттығу}

\pic4r{-12pt}{0.171}{figures/3.2.png}[\taskref{1}--\taskref{3} пункттеріне арналған соленоид]

Радиусы $r=\qty{1}{\cm}$ және ұзындығы $l=\qty{50}{\cm}$ цилиндрлік өзекшені қарастырайық, оның бүкіл ұзындығы бойынша $N=200$ орам сым біркелкі оралған. Сым бойымен $I=\qty{1}{\A}$ тұрақты ток ағып жатыр.
\begin{task}
    \item Шектік эффектілерді ескермей, соленоид ортасындағы магнит индукциясының $B$ модулін есептеңіз.
\end{task}

\begin{task}
    \item Өзекшедегі магнит ағынын $\Phi$-ді келесі түрде жазуға болады:
    \begin{equation}\label{task:3.1}
    \Phi = \frac{NI}{R_m},
    \end{equation}
    мұндағы $R_m$ --- \textit{магниттік кедергі} деп аталады, ал $\mathcal F=NI$ кейде \textit{магниттік қозғаушы күш} деп аталады. Алдыңғы пункттің аясында $R_m$-ді есептеңіз.
    \item Осындай соленоидтың $L$ индуктивтілігін сандық түрде есептеңіз.
\end{task}

\pic6r{0pt}{0.238}{figures/3.3.png}[\taskref{4} пунктіне; масштаб сақталмаған]

Пішіні күрделірек және иілген өзекшелер үшін магнит ағыны оның пішінін «қайталайды» және шашырау ағыны іс жүзінде түзілмейді. Бұған келесі пунктте көз жеткізуге болады:
\begin{task}
    \item \figref{2}-суретте көрсетілгендей ауа-темір жазық шекарасын қарастырайық: ауадағы магнит индукциясының $\vec B_0$ векторы нормальға $\alpha=\ang{10}$ бұрыш жасай бағытталған. Темірдегі магнит өрісінің $\vec B$ векторы нормальға $\beta$ бұрышымен көлбеуленген делік. $\beta$ неге тең және екі орта арасындағы магнит өрістері модульдерінің $B/B_0$ қатынасы қандай? Беттік өткізгіштік токтар жоқ.
\end{task}

\part{Саңылауы бар өзекше}

\pic9R{0pt}{0.226}{figures/3.4.png}[\taskref{5}--\taskref{7} пункттеріне; масштаб сақталмаған]

Қабырғасы $a=\qty{20}{\cm}$ және қимасының ауданы $S$ болатын жұмырланған квадрат пішінді өзекшені қарастырайық. Квадраттың бір жағына $N$ орам сым тығыз оралған, ол арқылы $I$ тогы өтеді, ал қарама-қарсы жағында қалыңдығы $d$ болатын саңылау кесілген (\figref{3}-сурет).

\begin{task}
    \item Бүкіл жүйенің $R_m=R_0$ магниттік кедергісін табыңыз, мұндағы $R_0$ \eqref{task:3.1} бойынша анықталады. Оны өзекшенің және ауа саңылауының магниттік кедергілерінің қосындысы ретінде жазуға болатынын көрсетіңіз.
\end{task}

\begin{task}
    \item $d$-ның қандай сандық мәнінде өзекшенің және ауа саңылауының магниттік кедергілері тең болады?
\end{task}
\newpage

Геометриялық өлшемдері \figref{3}-суреттегідей, \textit{бірақ ауа саңылауы жоқ} ($d=0$) квадрат өзекшенің қарама-қарсы жақтарына екі сым оралған: бірінші сымның $N_1$ орамы бар және ол қандай да бір $V(t)$ айнымалы кернеу желісіне қосылған; екінші сымның $N_2$ орамы бар және ол кедергісі $R$ резисторға қосылған. Шашырау ағыны жоқ.
\begin{task}
    \item Екі катушканың өзара индукция коэффициенті $M$-ді анықтаңыз. $V(t)$ қорек көзі жағынан көрінетін тізбектің $R'$ кедергісі неге тең? \textit{Кеңес 1:} $M$ коэффициенті бірінші катушканың екінші катушкамен ілініскен $\Phi_{12}$ магнит ағынының бірінші катушкадағы токқа қатынасы ретінде анықталады: $M=\frac{N_2\Phi_{12}}{I_2}=\frac{N_1\Phi_{21}}{I_1}$. \textit{Кеңес 2:} $R'$-ті табу үшін өзекшенің $R_m$ кедергісінің аздығын ескеріңіз.
\end{task}

\part{Электромагниттік аналогия}

\taskref{1}--\taskref{6} пункттері электрлік және магниттік тізбектер арасында аналогия жүргізуге итермелейді. Бұл аналогия \textit{есептің келесі пункттерін шешуге пайдалы болады}.

\begin{task}
    \item Төменде электрлік және магниттік шамалардың сәйкестік кестесі берілген: кестені жұмыс парағыңызға көшіріп, магниттік шамалардың тұсына оларға сәйкес келетін электрлік шамаларды және керісінше жазуыңыз керек. \textit{Егер сөзбен анықтамасын жазу қиын болса, міндетті түрде берілген шаманың стандартты белгісін көрсетіңіз.} Алғашқы екі жол сіз үшін толтырылған; оларды жұмыс парағына көшірудің қажеті жоқ.
    \begin{center}
    \begin{tabular}{|c|c|}\hline
         \textbf{Магниттік тізбек}& \textbf{Электрлік тізбек} \\\hline
         Магниттік қозғаушы күш $NI$& Электрлік қозғаушы күш $\mathcal E$\\\hline
         Магниттік кедергі $R_m$& Электрлік кедергі $R$\\\hline
         Магнит ағыны $\Phi$ &\\\hline
         &Электр өрісінің кернеулігі $\vec E$ \\\hline
         Абсолют магниттік өтімділік $\mu\mu_0$&\\\hline
    \end{tabular}
    \end{center}

    \item Магниттік қозғаушы күші $\mathcal F$ тұрақты болғанда, магниттік кедергісі $R_m$ болатын магниттік тізбек бөлігіндегі жинақталған $W_m$ магниттік энергия неге тең екенін көрсетіңіз.
\end{task}

\part{Дифференциалдық трансформатор}

Дифференциалдық трансформатордың қарапайым түрі \figref{4}-суретте көрсетілген. Өлшемдері $2l\times l$ ($l=\qty{20}{\cm}$) және ұзындық бірлігіне келетін магниттік кедергісі $\rho_m=\qty{4e5}{\per\H\per\m}$ болатын «8» пішінді өзекшенің екі қарама-қарсы шетіне орам сандары бірдей $N=10$ сымдар оралған. Электротехникада сол жақтағы сымды \textit{фаза} --- кернеу астындағы желі --- деп атайды және ол арқылы амплитудалық мәні $I_1=\qty{5}{\A}$ және жиілігі $f=\qty{50}{Гц}$ айнымалы ток өтеді; оң жақтағы сымды \textit{нөл} --- нөлдік жұмысшы өткізгіш --- деп атайды, ол арқылы қарама-қарсы бағытта амплитудалық мәні $I_2=\qty{4.97}{\A}$ және жиілігі $f=\qty{50}{Гц}$ айнымалы ток өтеді. $I_1$ мен $I_2$ арасында фазалық ығысу жоқ деп есептеңіз. Өзекшенің ортасында $N_0=1000$ орамы бар екінші реттік орам орналасқан, ол тізбектей жалғанған $R=\qty{200}{\ohm}$ резисторға және $M$ магниттік іске қосқышына (пускатель) қосылған. $M$-нің активті және реактивті кедергілері өте аз деп есептеңіз.
\begin{task}
    \item $M$ магниттік іске қосқышы арқылы өтетін $I_M$ амплитудалық ток күшін анықтаңыз.
\end{task}

\cpic{0.5}{figures/3.5.png}[\taskref{10} пунктіне арналған дифференциалдық трансформатор]

\part{Магниттік іске қосқыш}

\figref{5}-суретте көрсетілген магниттік іске қосқыштың (пускатель) жұмыс істеу принципін егжей-тегжейлі қарастырайық. Радиусы $=\qty{15}{\mm}$ және қалыңдығы аз П-тәрізді тірекке \textbf{(a)} радиусы $a=\qty{7}{\mm}$ темір өзекше \textbf{(c)} бекітілген, ал тірекке \textbf{(a)} бекітілген жалпы қатаңдығы $k=\qty{50}{\N\per\m}$ серіппелер \textbf{(e)} арқылы жылжымалы темір өзекше \textbf{(b)} жалғанған. Жылжымалы өзекшенің ішкі жағы (торец) мен \textbf{(c)} өзекшесінің жоғарғы жағы (торец) арасында бастапқы биіктігі $d=\qty{5}{\mm}$ болатын вертикаль ауа саңылауы қалады. Жылжымалы өзекшенің ішкі радиусы $a$-дан сәл үлкен, ал сыртқы радиусы $b$-дан сәл кіші (яғни \textbf{(a)} тірегі мен \textbf{(b)}, \textbf{(c)} өзекшелері арасындағы жұптық саңылаулар $d$-дан әлдеқайда аз). Өзекшелердің ұзындығы $b$-дан әлдеқайда үлкен. Тіректің айналасына сыртқы тізбекке (суретте көрсетілмеген) қосылған сым орамдары \textbf{(d)} тығыз оралған, бұл әсер етуші $\mathcal F$ магниттік қозғаушы күшін қамтамасыз етеді. \textbf{(b)} өзекшесінің жоғарғы жағына өткізгіш көпіршесі \textbf{(g)} мен екі клеммасы \textbf{(f)} бар бекіткіш орнатылған. \textbf{(f)} клеммаларынан $x_c=\qty{3}{\mm}$ қашықтықта сыртқы ажыратқышқа (суретте көрсетілмеген) қосылған \textbf{(h)} клеммалары орналасқан. Егер \textbf{(f)} клеммалары \textbf{(h)} клеммаларына тисе, ажыратқыш тізбегі тұйықталып, \taskref{10} пунктіндегі дифференциалдық трансформатордың қоректенуін өшіруге әкеледі. $\mathcal F=0$ болғанда серіппелердің \textbf{(e)} деформациясы жоқ; ауырлық күшін ескермеңіз; шеткі эффектілерді ескермеңіз.

\begin{task}
    \item $\mathcal F=\mathcal F_0$ қандай минималды мәнінде магниттік іске қосқыш тұйықталып, дифференциалдық трансформатор қорек көзінен ажыратылады? Егер барлығы $N=1000$ орам \textbf{(d)} болса, магниттік іске қосқыш орамындағы табалдырықты амплитудалық ток күші $I_M$ қандай болғанда сақтандырғыш іске қосылады? Нәтижені \taskref{10} пунктімен салыстырыңыз. Процестің квазистатикалық динамикасын қарастырыңыз.
\end{task}

Бұл есеп дифференциалдық трансформатордың тамаша қорғаныстық ажырату құрылғысы екенін көрсетеді --- сыртқы шашырау (адам, ылғал, жер және т.б. арқылы) салдарынан трансформаторға берілетін токтың шамалы өзгеруі \taskref{10} пунктіндегі $I_1\neq I_2$ жағдайына, ал бұл өз кезегінде қорғаныстық ажыратудың іске қосылуына әкеледі!

\cpic{0.48}{figures/3.6.png}[\taskref{11} пунктіне арналған магниттік іске қосқыш]
